%% GENERATED by tools/gen_latex.py -- edit content/ch01.py instead.
%% Build:  latexmk -pdf ch01.tex     (or: pdflatex ch01.tex, twice)
\documentclass[aspectratio=169,11pt,t]{beamer}
\usepackage{itbpro}

\coursetitle{Apa Itu Deep Learning?}
\coursesubtitle{Meletakkan AI, machine learning, dan deep learning pada tempatnya masing-masing -- lalu memisahkan capaian nyata dari gembar-gembor.}
\coursekicker{Chapter 1}
\coursesource{Chollet \& Watson, Deep Learning with Python 3e -- bab 1}
\courseduration{90 menit}
\coursepresenter{Prof. Bambang Riyanto Trilaksono}{Pengajar Utama}
\coursebrand{AI for Professional \textperiodcentered\ ITB \textperiodcentered\ Ch 1}

\title{Apa Itu Deep Learning?}
\author{Prof. Bambang Riyanto Trilaksono}
\date{}

\begin{document}
\covertitle

\framekicker{Learning outcomes}
\begin{frame}[shrink=25]{Tujuan Sesi}
\leadpar{Setelah sesi ini, peserta dapat:}
\begin{isteps}
\item Menempatkan \textbf{AI, machine learning, dan deep learning} dalam hubungan yang benar -- yang mana memuat yang mana, dan sejak kapan.
\item Menjelaskan \textbf{pembalikan paradigma}: pemrograman klasik menghasilkan jawaban, machine learning menghasilkan aturan.
\item Menyebut \textbf{tiga bahan wajib} setiap sistem machine learning dan menunjukkan apa yang rusak bila salah satunya hilang.
\item Menerangkan cara kerja deep learning lewat \textbf{bobot, fungsi rugi, dan backpropagation} tanpa menurunkan satu pun rumus.
\item Membedakan \textbf{capaian yang sudah terbukti} dari \textbf{klaim yang belum}, dan menyebut apa yang memicu dua musim dingin AI sebelumnya.
\end{isteps}

\vskip 6pt
\vskip 4pt
\reslink{Buku}{Bab 1 --- teks penuh}{https://deeplearningwithpython.io/chapters/chapter01\_what-is-deep-learning}
\reslink{Notebook}{01\_paradigma\_ml.ipynb}{../../notebooks/ch01/01\_paradigma\_ml.ipynb}
\vskip 2pt
\end{frame}

\framekicker{Peta bab}
\begin{frame}[shrink=25]{Ke mana bab ini membawa kita}
\leadpar{Bab pembuka tidak mengajarkan satu baris kode pun. Tugasnya lebih mendasar: memastikan kita memakai kata yang sama untuk hal yang sama.}
\begin{minipage}[t]{0.3253\textwidth}
\begin{icard}
\cardh{Definisi dan sejarah}
Tiga lingkaran bersarang: AI, ML, DL. Ditambah satu cabang yang \textbf{bukan} ML sama sekali -- symbolic AI.
\cardtag{bag. 1-3}
\end{icard}
\end{minipage}%
\hfill
\begin{minipage}[t]{0.3253\textwidth}
\begin{icard}
\cardh{Cara kerjanya}
Representasi, bobot, fungsi rugi, backpropagation. Chollet menerangkannya dengan \textbf{tiga gambar}, bukan rumus.
\cardtag{bag. 4-6}
\end{icard}
\end{minipage}%
\hfill
\begin{minipage}[t]{0.3253\textwidth}
\begin{icard}
\cardh{Capaian vs gembar-gembor}
Apa yang \textbf{sudah} dikerjakan deep learning, dan mengapa dua kali sebelumnya musim panas AI berubah jadi musim dingin.
\cardtag{bag. 7-12}
\end{icard}
\end{minipage}%
\begin{iband}[signal]
Untuk kelas ini bab 1 punya beban tambahan: ia yang menetapkan \textbf{kosakata bersama} yang dipakai sampai bab 20 dan sampai topik agentic AI di akhir kursus.
\end{iband}
\end{frame}
\note{Buka dengan pertanyaan ke peserta: siapa yang bisa menjelaskan beda AI dan machine learning dalam satu kalimat? Jawaban yang beragam di ruangan justru bahan yang bagus untuk masuk ke slide berikutnya.}

\coursesection{01}{AI, machine learning, deep learning}{Tiga istilah yang di media dipakai bergantian, padahal bersarang.}

\framekicker{Bagian 1.1}
\begin{frame}[shrink=25]{Tiga lingkaran yang bersarang, bukan tiga sinonim}
\begin{center}
\adjustbox{max width=\linewidth,max totalheight=0.52\textheight}{%

\begin{tikzpicture}[font=\sffamily]
  \node[draw=itbbluelt!70, fill=itbbluelt!8, rounded corners=5pt,
        minimum width=7.4cm, minimum height=3.4cm, anchor=north west] (ai) at (0,0) {};
  \node[anchor=north west, font=\bfseries\small, text=ink] at ($(ai.north west)+(0.28,-0.22)$)
    {Artificial Intelligence};
  \node[anchor=north west, font=\tiny, text=ink3] at ($(ai.north west)+(0.28,-0.62)$)
    {sejak 1950-an --- ``mengotomasi tugas intelektual''};

  \node[draw=signal!70, fill=signal!8, rounded corners=5pt,
        minimum width=6.5cm, minimum height=2.2cm, anchor=north west] (ml) at (0.45,-1.05) {};
  \node[anchor=north west, font=\bfseries\small, text=ink] at ($(ml.north west)+(0.26,-0.2)$)
    {Machine Learning};
  \node[anchor=north west, font=\tiny, text=ink3] at ($(ml.north west)+(0.26,-0.58)$)
    {sejak 1990-an --- aturan dipelajari dari data};

  \node[draw=violet!70, fill=violet!10, rounded corners=5pt,
        minimum width=5.6cm, minimum height=1.0cm, anchor=north west] (dl) at (0.9,-2.0) {};
  \node[anchor=north west, font=\bfseries\small, text=ink] at ($(dl.north west)+(0.24,-0.18)$)
    {Deep Learning};
  \node[anchor=north west, font=\tiny, text=ink3] at ($(dl.north west)+(0.24,-0.54)$)
    {lapisan representasi bertingkat};

  \node[draw=rule, fill=papertint, rounded corners=5pt,
        minimum width=3.4cm, minimum height=2.1cm, anchor=north west] (sym) at (8.0,-0.6) {};
  \node[anchor=north west, font=\bfseries\small, text=ink] at ($(sym.north west)+(0.24,-0.2)$)
    {Symbolic AI};
  \node[anchor=north west, font=\tiny, text=ink3, align=left] at ($(sym.north west)+(0.24,-0.62)$)
    {1950-an -- 1980-an\\aturan ditulis tangan\\unggul di catur,\\gagal di persoalan kabur};
  \draw[rule, dashed, line width=0.6pt] (7.45,-1.7) -- (7.95,-1.7);
\end{tikzpicture}

}
\end{center}
\vskip 3pt{\color{ink3}\fontsize{7.4}{9.4}\selectfont Gambar 1.1 -- deep learning adalah bagian dari machine learning, yang adalah bagian dari AI. Symbolic AI ada di dalam AI tetapi di luar machine learning.\par}
\begin{ibullets}
\item \textbf{AI} lahir 1950-an. Definisi kerjanya: \textit{upaya mengotomasi tugas intelektual yang biasanya dikerjakan manusia}. Cakupannya jauh lebih luas dari ML.
\item \textbf{Symbolic AI} mendominasi 1950-an sampai 1980-an: aturan ditulis tangan oleh pemrogram. Cukup untuk catur, \hilite{runtuh} di persoalan kabur seperti mengenali gambar atau bahasa.
\item \textbf{Machine learning} naik sejak 1990-an justru karena persoalan kabur itulah yang tersisa.
\end{ibullets}
\end{frame}
\note{Tekankan: kegagalan symbolic AI bukan kegagalan gagasan AI. Ia gagal pada satu kelas persoalan tertentu -- yang aturannya tidak bisa dituliskan manusia karena manusia sendiri tidak sadar memakainya.}

\framekicker{Bagian 1.1}
\begin{frame}[shrink=25]{Kutipan yang membingkai seluruh bidang}
\begin{iquote}
\quotetext{The Analytical Engine has no pretensions whatever to originate anything. It can do whatever we know how to order it to perform.}\quotecite{Ada Lovelace, 1843 -- catatan atas mesin Charles Babbage}
\end{iquote}
\begin{iquote}
\quotetext{Every aspect of learning or any other feature of intelligence can in principle be so precisely described that a machine can be made to simulate it.}\quotecite{John McCarthy dkk., proposal lokakarya Dartmouth, 1956}
\end{iquote}
\begin{iband}[amber]
Dua kutipan ini berselisih 113 tahun dan berselisih pendapat. Pertanyaan Lovelace -- \textit{bisakah mesin melampaui perintah kita?} -- belum tuntas sampai hari ini; \hilite{machine learning adalah jawaban parsialnya}, sebab mesin memang menghasilkan sesuatu yang tidak kita tuliskan: aturannya.
\end{iband}
\end{frame}
\note{Kutipan Lovelace dipakai Chollet sejak edisi pertama sebagai pembuka. Sambungkan ke pertanyaan yang akan diajukan auditor mana pun: kalau aturannya tidak kita tulis, siapa yang bertanggung jawab atas aturan itu?}

\coursesection{02}{Pembalikan paradigma}{Yang tadinya keluaran, kini menjadi masukan.}

\framekicker{Bagian 1.1.2}
\begin{frame}[shrink=25]{Machine learning membalik arah pemrograman}
\begin{center}
\adjustbox{max width=\linewidth,max totalheight=0.52\textheight}{%

\begin{tikzpicture}[font=\sffamily\tiny,
  bx/.style={draw=rule, fill=papertint, rounded corners=3pt, minimum width=1.9cm,
             minimum height=0.6cm, text=ink2},
  ax/.style={draw=signal!60, fill=signal!10, rounded corners=3pt, minimum width=1.8cm,
             minimum height=0.6cm, text=ink},
  vx/.style={draw=violet!70, fill=violet!12, rounded corners=3pt, minimum width=1.8cm,
             minimum height=0.6cm, text=ink},
  ar/.style={-{Stealth[length=4pt]}, signal, line width=0.7pt}]

  \node[font=\bfseries\scriptsize, text=ink, anchor=west] at (0,1.55) {Pemrograman klasik};
  \node[bx] (r1) at (1,1.0) {Aturan};
  \node[bx] (d1) at (1,0.2) {Data};
  \node[ax] (p1) at (3.6,0.6) {Program};
  \node[bx] (a1) at (6.1,0.6) {Jawaban};
  \draw[ar] (r1) -- (p1); \draw[ar] (d1) -- (p1); \draw[ar] (p1) -- (a1);

  \draw[rule, line width=0.5pt] (0,-0.4) -- (8.6,-0.4);

  \node[font=\bfseries\scriptsize, text=ink, anchor=west] at (0,-0.95) {Machine learning};
  \node[bx] (d2) at (1,-1.4) {Data};
  \node[bx] (a2) at (1,-2.2) {Jawaban};
  \node[ax] (p2) at (3.6,-1.8) {Pelatihan};
  \node[vx] (r2) at (6.1,-1.8) {Aturan};
  \draw[ar] (d2) -- (p2); \draw[ar] (a2) -- (p2); \draw[ar] (p2) -- (r2);
  \node[text=ink3, anchor=west, align=left] at (7.2,-1.8)
    {yang tadinya keluaran\\kini menjadi masukan};
\end{tikzpicture}

}
\end{center}
\vskip 3pt{\color{ink3}\fontsize{7.4}{9.4}\selectfont Gambar 1.2 -- pemrograman klasik disuapi aturan dan data lalu mengeluarkan jawaban; machine learning disuapi data dan jawaban lalu mengeluarkan aturan.\par}
\begin{iband}[signal]
Konsekuensinya langsung dan sering diremehkan: sistem ML \hilite{dilatih, bukan diprogram}. Kalau tidak ada contoh yang sudah berjawab, tidak ada yang bisa dilatih -- seberapa pun bagus arsitekturnya.
\end{iband}
\end{frame}
\note{Di sini biasanya muncul pertanyaan 'kalau begitu datanya dari mana?'. Jawab singkat saja, bab 6 membahasnya penuh sebagai universal workflow.}

\framekicker{Bagian 1.1.2}
\begin{frame}[shrink=25]{Tiga bahan wajib -- hilang satu, bukan machine learning}
\begin{minipage}[t]{0.3253\textwidth}
\begin{icardaccent}
\cardh{1 $\cdot$ Titik data masukan}
Berkas suara, gambar, baris transaksi. Inilah yang akan ditransformasi.
\end{icardaccent}
\end{minipage}%
\hfill
\begin{minipage}[t]{0.3253\textwidth}
\begin{icardaccent}
\cardh{2 $\cdot$ Contoh keluaran yang diharapkan}
Transkrip untuk suara, label untuk gambar, penanda \textit{fraud} untuk transaksi.
\end{icardaccent}
\end{minipage}%
\hfill
\begin{minipage}[t]{0.3253\textwidth}
\begin{icardaccent}
\cardh{3 $\cdot$ Cara mengukur mutu}
Jarak antara tebakan dan jawaban benar. Inilah \textbf{sinyal umpan balik} yang menyetel algoritmanya.
\end{icardaccent}
\end{minipage}%
{\fontsize{9.4}{12.6}\selectfont\color{fg2}Ketiganya adalah syarat minimum, bukan daftar keinginan. Chollet menyebut inti persoalannya sebagai \textit{meaningfully transform data} -- mempelajari \textbf{representasi} yang membuat tugas jadi lebih mudah.\par}\vskip 4pt
\begin{itable}{>{\raggedright\arraybackslash}p{0.2444\textwidth}>{\raggedright\arraybackslash}p{0.3196\textwidth}>{\raggedright\arraybackslash}p{0.376\textwidth}}
\thead{Bahan yang hilang} & \thead{Yang sebenarnya Anda punya} & \thead{Akibatnya} \\ \midrule
Contoh keluaran & Data mentah tanpa label & Bukan supervised learning. Pilihannya: pelabelan, atau pindah ke unsupervised / self-supervised. \\
Ukuran mutu & Data dan label, tanpa metrik yang disepakati & Model bisa dilatih tetapi tidak bisa dinilai -- dan ini kegagalan proyek yang paling sering di dunia nyata. \\
Data masukan yang mewakili & Contoh yang tidak menyerupai kondisi produksi & Model berhasil di laptop, gagal saat dipakai. Bab 5 dan 6 menamainya secara teknis. \\
\end{itable}
\end{frame}
\note{Minta peserta memikirkan satu kasus pemakaian dari pekerjaannya sendiri, lalu uji ketiga bahan itu. Yang paling sering hilang adalah bahan ketiga -- kesepakatan tentang apa yang disebut 'benar'.}

\framekicker{Bagian 1.1.3}
\begin{frame}[shrink=25]{Representasi: gagasan yang menyatukan seluruh buku}
\begin{minipage}[t]{0.485\textwidth}
{\fontsize{9.4}{12.6}\selectfont\color{fg2}\textbf{Representasi} adalah cara lain menyandikan data yang sama. Foto yang sama bisa disandikan sebagai RGB atau HSV -- isinya identik, tetapi tugas yang mudah pada masing-masing berbeda.\par}\vskip 4pt
\begin{ibullets}
\item \textit{Pilih semua piksel merah} -- mudah di \textbf{RGB}.
\item \textit{Buat gambar kurang jenuh} -- mudah di \textbf{HSV}.
\item Persoalan yang sama; yang berubah hanya sumbunya.
\end{ibullets}
\begin{iband}[signal]
Karena itu Chollet merumuskan machine learning sebagai \hilite{pencarian representasi yang membuat tugas jadi mudah}.
\end{iband}

\end{minipage}%
\hfill
\begin{minipage}[t]{0.485\textwidth}
{\fontsize{9.4}{12.6}\selectfont\color{fg2}Contoh klasik bab ini: titik hitam dan putih yang bercampur pada sumbu asal (gambar 1.3). Tidak ada aturan sederhana yang memisahkannya.\par}\vskip 4pt
{\fontsize{9.4}{12.6}\selectfont\color{fg2}Ubah sumbunya -- pindahkan titik asal, putar -- dan aturannya menjadi satu kalimat (gambar 1.4): \textit{hitam bila x > 0}.\par}\vskip 4pt
\begin{iband}[amber]
Tidak ada model yang jadi lebih pintar. Yang berubah hanya representasinya. Inilah pekerjaan yang dulu disebut \textbf{feature engineering} dan dikerjakan manusia.
\end{iband}

\end{minipage}%
\end{frame}
\note{Kalau ada waktu, gambar manual di papan: sebar titik, lalu putar sumbu. Peraga fisik ini jauh lebih nempel daripada slide.}

\framekicker{Bagian 1.4}
\begin{frame}[shrink=25]{Ruang hipotesis -- dan mengapa aturan tulis-tangan runtuh}
{\fontsize{9.4}{12.6}\selectfont\color{fg2}Perubahan sumbu tadi kita rancang \textbf{dengan tangan}. Itu sanggup untuk persoalan sesederhana itu. Tetapi bisakah Anda menuliskan transformasi citra yang menjelaskan beda 6 dan 8, atau 1 dan 7, \hilite{untuk segala macam tulisan tangan}?\par}\vskip 4pt
\begin{minipage}[t]{0.494\textwidth}
\begin{icardwarn}
\cardh{Bisa -- sampai batas tertentu}
Aturan seperti \textit{menghitung banyaknya gelung tertutup}, atau histogram piksel tegak dan mendatar, lumayan membedakan angka tulisan tangan.
\end{icardwarn}
\end{minipage}%
\hfill
\begin{minipage}[t]{0.494\textwidth}
\begin{icardbad}
\cardh{Tetapi rapuh dan menyiksa}
Tiap kali muncul contoh tulisan baru yang mematahkan aturan yang sudah disusun rapi, Anda harus menambah transformasi dan aturan baru -- \textbf{sambil memperhitungkan interaksinya dengan semua aturan sebelumnya}.
\end{icardbad}
\end{minipage}%
\begin{iband}[signal]
Algoritma machine learning \textbf{tidak kreatif} dalam menemukan transformasi ini. Ia sekadar menelusuri sehimpunan operasi yang sudah ditetapkan lebih dulu -- himpunan itulah yang disebut \hilite{ruang hipotesis (hypothesis space)}. Pada contoh 2D tadi, ruang hipotesisnya adalah ruang semua perubahan sumbu yang mungkin.
\end{iband}
\begin{iquote}
\quotetext{Itulah machine learning, secara ringkas: \textbf{mencari representasi dan aturan yang berguna atas suatu data masukan, di dalam ruang kemungkinan yang sudah ditetapkan, dengan tuntunan sebuah sinyal umpan balik.}}\quotecite{Chollet \& Watson, bab 1.4}
\end{iquote}
{\fontsize{9.4}{12.6}\selectfont\color{fg2}Gagasan sesederhana itu menyelesaikan rentang tugas intelektual yang luar biasa lebar -- dari kemudi otonom sampai penjawaban pertanyaan dalam bahasa alami.\par}\vskip 4pt
\end{frame}
\note{Istilah 'ruang hipotesis' akan kembali di bab 3 saat topologi model dibahas: memilih arsitektur = mempersempit ruang hipotesis.}

\coursesection{03}{Yang 'dalam' pada deep learning}{Bukan pemahaman yang lebih dalam. Hanya lapisan yang lebih banyak.}

\framekicker{Bagian 1.1.4}
\begin{frame}[shrink=25]{'Deep' menunjuk pada jumlah lapis, bukan kedalaman makna}
\begin{iquote}
\quotetext{The "deep" in "deep learning" isn't a reference to any kind of deeper understanding -- rather, it stands for this idea of successive layers of representations.}\quotecite{Chollet \& Watson, bab 1}
\end{iquote}
\begin{center}
\adjustbox{max width=\linewidth,max totalheight=0.52\textheight}{%

\begin{tikzpicture}[font=\sffamily\tiny,
  lay/.style={draw=signal!60, fill=signal!9, rounded corners=3pt,
              minimum width=1.25cm, minimum height=1.5cm, text=ink, align=center},
  io/.style={draw=rule, fill=papertint, rounded corners=3pt,
             minimum width=1.25cm, minimum height=1.2cm, text=ink2, align=center},
  ar/.style={-{Stealth[length=4pt]}, signal, line width=0.7pt}]
  \node[io]  (in) at (0,0)   {masukan\\$28\times28$};
  \node[lay] (l1) at (1.9,0) {lapis 1\\\ttfamily tepi};
  \node[lay] (l2) at (3.5,0) {lapis 2\\\ttfamily sudut};
  \node[lay] (l3) at (5.1,0) {lapis 3\\\ttfamily bagian};
  \node[lay] (l4) at (6.7,0) {lapis 4\\\ttfamily angka};
  \node[draw=lime!60, fill=limebr!12, rounded corners=3pt, minimum width=1.6cm,
        minimum height=1.2cm, text=ink, align=center] (out) at (8.7,0) {keluaran\\\ttfamily P(0..9)};
  \draw[ar] (in) -- (l1); \draw[ar] (l1) -- (l2); \draw[ar] (l2) -- (l3);
  \draw[ar] (l3) -- (l4); \draw[ar] (l4) -- (out);
  \node[text=ink3, anchor=west] at (-0.7,-1.25)
    {setiap lapis mengubah representasi menjadi bentuk yang sedikit lebih berguna};
\end{tikzpicture}

}
\end{center}
\vskip 3pt{\color{ink3}\fontsize{7.4}{9.4}\selectfont Gambar 1.5-1.6 -- jaringan empat lapis untuk klasifikasi angka. Representasi antara semakin jauh dari piksel dan semakin dekat ke jawaban.\par}
\begin{ibullets}
\item Banyaknya lapis = \textbf{kedalaman} model. Jaringan modern punya puluhan sampai ratusan lapis.
\item Lawannya, \textbf{shallow learning}, hanya satu atau dua lapis representasi.
\item Semua lapis dipelajari \hilite{sekaligus, otomatis} dari data pelatihan -- bukan dirancang satu per satu oleh manusia.
\item Chollet menyebut nama yang sebenarnya lebih tepat untuk bidang ini: \textit{layered representations learning} atau \textit{hierarchical representations learning}. Istilah \textbf{deep} menang karena sejarah, bukan karena akurasi.
\end{ibullets}
\end{frame}
\note{Koreksi salah paham yang paling umum di ruangan: 'deep' sering dikira berarti mesin memahami lebih dalam. Bukan. Ia sekadar keterangan arsitektur.}

\framekicker{Bagian 1.1.5 -- tiga gambar}
\begin{frame}[shrink=25]{Cara kerja deep learning, tanpa satu pun rumus}
\begin{center}
\adjustbox{max width=\linewidth,max totalheight=0.52\textheight}{%

\begin{tikzpicture}[font=\sffamily\tiny,
  bx/.style={draw=rule, fill=papertint, rounded corners=3pt, minimum width=1.7cm,
             minimum height=0.62cm, text=ink2, align=center},
  ar/.style={-{Stealth[length=4pt]}, signal, line width=0.7pt},
  fb/.style={-{Stealth[length=4pt]}, amberbr, line width=0.8pt}]
  \node[bx] (x) at (0,0) {Masukan $X$};
  \node[draw=itbbluelt!70, fill=itbbluelt!12, rounded corners=4pt, minimum width=2.1cm,
        minimum height=1.7cm, text=ink, align=center] (lay) at (2.5,0)
        {Lapis\\(transformasi data)};
  \node[draw=amber!70, fill=amberbr!16, rounded corners=3pt, minimum width=1.6cm,
        minimum height=0.42cm, text=ink, font=\ttfamily\tiny] (w) at (2.5,-0.62) {Bobot $W$};
  \node[bx] (yp) at (5.2,0) {Prediksi $\hat{Y}$};
  \node[bx] (y)  at (5.2,-1.5) {Target $Y$};
  \node[draw=rose!70, fill=rosebr!14, rounded corners=4pt, minimum width=1.9cm,
        minimum height=0.62cm, text=ink] (loss) at (7.9,-0.75) {Fungsi rugi};
  \node[draw=lime!70, fill=limebr!16, rounded corners=4pt, minimum width=1.9cm,
        minimum height=0.62cm, text=ink] (opt) at (7.9,1.15) {Optimalisasi};

  \draw[ar] (x) -- (lay); \draw[ar] (lay) -- (yp);
  \draw[ar] (yp) -- (loss); \draw[ar] (y) -- (loss);
  \draw[fb] (loss) -- (opt);
  \draw[fb, dashed] (opt) -| (2.5,0.9) -- (lay.north);
  \node[text=amber, anchor=south] at (4.6,1.3) {perbarui bobot};
\end{tikzpicture}

}
\end{center}
\vskip 3pt{\color{ink3}\fontsize{7.4}{9.4}\selectfont Gambar 1.7-1.9 digabung: bobot memarametrikan lapis, fungsi rugi mengukur jaraknya ke target, optimalisasi memakai skor itu untuk menggeser bobot.\par}
\begin{isteps}
\item \textbf{Bobot memarametrikan transformasi} (gambar 1.7). Apa yang dikerjakan sebuah lapis tersimpan di bobotnya. Belajar = menemukan nilai bobot yang benar. Jaringan besar punya puluhan juta parameter.
\item \textbf{Fungsi rugi mengukur seberapa jauh melesetnya} (gambar 1.8). Ia menghitung jarak antara prediksi dan target. Disebut juga \textit{objective} atau \textit{cost function}.
\item \textbf{Backpropagation memakai skor itu sebagai umpan balik} (gambar 1.9): geser bobot ke arah yang menurunkan rugi, ulangi ribuan kali. Inilah \textbf{lingkar pelatihan}.
\end{isteps}
\begin{iband}[signal]
Bobot mula-mula diisi \hilite{acak}, jadi keluaran awal tentu ngawur dan ruginya tinggi. Yang membuatnya bekerja adalah pengulangan lingkar itu, bukan tebakan awal yang bagus.
\end{iband}
\end{frame}
\note{Ini slide paling penting di bab 1. Bab 2 akan membongkar setiap kotak di diagram ini menjadi tensor dan turunan; cukup pastikan bentuk lingkarnya tertanam dulu.}

\framekicker{Peraga 1 dari 2}
\begin{frame}[shrink=25]{Lingkar pelatihan itu, dalam sepuluh baris}
{\fontsize{9.4}{12.6}\selectfont\color{fg2}Bab 1 memang tidak memuat kode. Tetapi seluruh diagram di slide sebelumnya muat dalam potongan sekecil ini.\par}\vskip 4pt
\icode{python}{peraga kelas --- bukan listing buku}{listings/ch01-01.py}
\end{frame}
\note{Jalankan langsung di depan kelas kalau memungkinkan. Tiga komentar bernomor di kode itu menunjuk persis ke tiga gambar pada slide sebelumnya -- tunjuk bolak-balik antara keduanya.}

\framekicker{Peraga 2 dari 2}
\begin{frame}[shrink=25]{Yang keluar dari lingkar itu: aturannya sendiri}
\ioutput{listings/ch01-02.txt}
\begin{iband}[signal]
Bobot bergerak dari acak ke \hilite{W $\approx$ 2, b $\approx$ 1} -- hukum yang memang membangkitkan datanya, dan yang tidak pernah kita tuliskan. Itulah 'aturan sebagai keluaran' pada gambar 1.2.
\end{iband}
\begin{ibullets}
\item Tidak ada satu baris pun yang memberi tahu program bahwa jawabannya 2 dan 1.
\item Yang diberikan hanya \textbf{data, jawaban, dan ukuran rugi} -- tiga bahan wajib itu.
\item Bab 2 mengganti NumPy sebaris ini dengan tensor dan turunan otomatis; bentuk lingkarnya \hilite{tidak berubah}.
\end{ibullets}
\vskip 4pt
\reslink{NOTEBOOK}{01\_paradigma\_ml.ipynb}{../../notebooks/ch01/01\_paradigma\_ml.ipynb}
\reslink{BAB PENUH}{deeplearningwithpython.io}{https://deeplearningwithpython.io/chapters/chapter01\_what-is-deep-learning}
\vskip 2pt
\end{frame}
\note{Yang mau dilihat peserta bukan angkanya, melainkan bahwa loss turun sendiri tanpa ada yang memberitahu jawabannya.}

\coursesection{04}{Mengapa deep learning menang}{Tiga sifat, dan satu perubahan yang membuat pesaingnya tertinggal.}

\framekicker{Bagian 1.1.6}
\begin{frame}[shrink=25]{Apa yang membuatnya berbeda dari shallow learning}
\begin{minipage}[t]{0.3253\textwidth}
\begin{icardgood}
\cardh{Kesederhanaan}
Feature engineering dikerjakan sendiri oleh model. Pipeline bertingkat digantikan \textbf{satu model ujung-ke-ujung}.
\cardtag{simplicity}
\end{icardgood}
\end{minipage}%
\hfill
\begin{minipage}[t]{0.3253\textwidth}
\begin{icardgood}
\cardh{Skalabilitas}
Cocok diparalelkan di GPU dan dilatih per \textit{batch}, jadi ukuran data \hilite{tidak lagi jadi batas atas}.
\cardtag{scalability}
\end{icardgood}
\end{minipage}%
\hfill
\begin{minipage}[t]{0.3253\textwidth}
\begin{icardgood}
\cardh{Keluwesan \& pakai ulang}
Model bisa dilatih lanjut dengan data baru tanpa mulai dari nol; \textit{foundation model} dipindah ke tugas lain dengan sedikit pelatihan tambahan.
\cardtag{reusability}
\end{icardgood}
\end{minipage}%
{\fontsize{9.4}{12.6}\selectfont\color{fg2}Perbedaan pokoknya satu: pendekatan dangkal memaksa manusia merancang lapisan representasi dulu, lalu model bekerja di atasnya. Deep learning \hilite{mempelajari semua lapisan itu sekaligus, dalam satu tarikan}.\par}\vskip 4pt
\begin{iband}[amber]
Ini juga sumber keluhan auditor: yang membuatnya kuat -- representasi yang ditemukan sendiri -- persis yang membuatnya sulit diterangkan. Bab 10 kembali ke soal ini untuk kasus citra.
\end{iband}
\end{frame}
\note{Sambungkan ke topik kepatuhan yang akan dibahas di sesi lain: keterjelasan model bukan sekadar keinginan akademik di sektor diatur ketat, ia persyaratan.}

\framekicker{Bagian 1.2}
\begin{frame}[shrink=25]{Zaman AI generatif}
\begin{minipage}[t]{0.575\textwidth}
{\fontsize{9.4}{12.6}\selectfont\color{fg2}Mekanismenya: model menyusun ulang teks atau gambar dari masukannya sendiri. Targetnya diambil \textbf{dari masukan itu juga} -- inilah \textit{self-supervised learning}, dan itulah sebabnya ia lepas dari batas ketersediaan label.\par}\vskip 4pt
\begin{ibullets}
\item \textit{Foundation model} dengan \textbf{ratusan miliar parameter}, dilatih di atas data lebih dari \textbf{satu petabyte}.
\item Berperilaku seperti \hilite{basis data kabur atas pengetahuan manusia}.
\item Persoalan baru diselesaikan lewat \textbf{prompting}, tanpa pemrograman khusus per tugas.
\item Masuk kesadaran umum pada \textbf{2022}, tetapi percobaan pembangkitan teks sudah ada sejak \textbf{1990-an}.
\end{ibullets}

\end{minipage}%
\hfill
\begin{minipage}[t]{0.395\textwidth}
\begin{minipage}[t]{1.0\textwidth}
\istat{10$^{1}$$^{1}$}{orde parameter pada foundation model terkini}
\end{minipage}%


\vskip 4pt

\begin{minipage}[t]{1.0\textwidth}
\istat{> 1 PB}{orde data pelatihan}
\end{minipage}%


\vskip 4pt

\begin{minipage}[t]{1.0\textwidth}
\istat{2022}{tahun ia menjadi perbincangan umum}
\end{minipage}%

\end{minipage}%
\begin{iband}[signal]
Untuk kursus ini, bab 15-17 membongkar mesinnya (Transformer, pembangkitan teks, pembangkitan gambar), dan topik LLM serta agentic AI memakainya sebagai bahan bangunan.
\end{iband}
\end{frame}
\note{Tahan dulu pertanyaan tentang RAG dan fine-tuning; itu topik 4. Di sini cukup tegaskan bahwa self-supervised-lah yang membuat skala sebesar itu mungkin.}

\coursesection{05}{Capaian, dan gembar-gembor}{Yang sudah terbukti, dan mengapa dua kali sebelumnya kita salah menaksir.}

\framekicker{Bagian 1.3}
\begin{frame}[shrink=25]{Yang sudah benar-benar dikerjakan deep learning}
\begin{minipage}[t]{0.241\textwidth}
\begin{icard}
\cardh{Percakapan}
ChatGPT, Gemini, Claude.
\end{icard}
\end{minipage}%
\hfill
\begin{minipage}[t]{0.241\textwidth}
\begin{icard}
\cardh{Pembangkitan kode}
GitHub Copilot dan sejenisnya.
\end{icard}
\end{minipage}%
\hfill
\begin{minipage}[t]{0.241\textwidth}
\begin{icard}
\cardh{Gambar fotorealistik}
Dari teks ke citra.
\end{icard}
\end{minipage}%
\hfill
\begin{minipage}[t]{0.241\textwidth}
\begin{icard}
\cardh{Setara manusia}
Klasifikasi citra, transkripsi suara, pengenalan tulisan tangan.
\end{icard}
\end{minipage}%


\vskip 4pt

\begin{minipage}[t]{0.241\textwidth}
\begin{icard}
\cardh{Terjemahan \& TTS}
Naik tajam mutunya.
\end{icard}
\end{minipage}%
\hfill
\begin{minipage}[t]{0.241\textwidth}
\begin{icard}
\cardh{Kemudi otonom}
Beroperasi di Phoenix, San Francisco, Los Angeles, Austin (2025).
\end{icard}
\end{minipage}%
\hfill
\begin{minipage}[t]{0.241\textwidth}
\begin{icard}
\cardh{Melampaui manusia}
Go, catur, poker.
\end{icard}
\end{minipage}%
\hfill
\begin{minipage}[t]{0.241\textwidth}
\begin{icard}
\cardh{Struktur protein}
AlphaFold.
\end{icard}
\end{minipage}%
\begin{minipage}[t]{0.3253\textwidth}
\begin{icard}
\cardh{Sistem perekomendasi}
YouTube, Netflix, Spotify -- yang paling banyak dipakai orang setiap hari tanpa menyadarinya.
\end{icard}
\end{minipage}%
\hfill
\begin{minipage}[t]{0.3253\textwidth}
\begin{icard}
\cardh{Naskah kuno}
Puluhan ribu manuskrip di \textbf{Vatican Secret Archive} ditranskripsi otomatis.
\end{icard}
\end{minipage}%
\hfill
\begin{minipage}[t]{0.3253\textwidth}
\begin{icard}
\cardh{Penyakit tanaman}
Dideteksi dan digolongkan langsung di lahan, cukup dengan ponsel biasa.
\end{icard}
\end{minipage}%


\vskip 4pt

\begin{minipage}[t]{0.3253\textwidth}
\begin{icard}
\cardh{Citra medis}
Mendampingi onkolog dan radiolog menafsirkan hasil pencitraan.
\end{icard}
\end{minipage}%
\hfill
\begin{minipage}[t]{0.3253\textwidth}
\begin{icard}
\cardh{Bencana alam}
Meramalkan banjir, badai, bahkan gempa bumi.
\end{icard}
\end{minipage}%
\hfill
\begin{minipage}[t]{0.3253\textwidth}
\begin{icard}
\cardh{Struktur protein}
AlphaFold, dengan ketepatan yang belum pernah ada sebelumnya.
\end{icard}
\end{minipage}%
\begin{iband}[signal]
Daftar ini \hilite{bukan janji}; semuanya sudah berjalan. Itulah yang membedakannya dari slide berikutnya.
\end{iband}
\end{frame}
\note{Sengaja ditaruh sebelum slide gembar-gembor. Urutannya penting: akui dulu capaiannya, baru kritik taksirannya. Kalau dibalik, terdengar seperti penyangkalan.}

\framekicker{Bagian 1.9}
\begin{frame}[shrink=25]{Tiga gelombang dalam satu dasawarsa}
\begin{minipage}[t]{0.3253\textwidth}
\begin{icardaccent}
\cardh{2013 -- 2017 $\cdot$ persepsi}
Hasil yang mencengangkan pada tugas persepsi: klasifikasi citra, transkripsi suara, pengenalan tulisan tangan.
\end{icardaccent}
\end{minipage}%
\hfill
\begin{minipage}[t]{0.3253\textwidth}
\begin{icardaccent}
\cardh{2017 -- 2022 $\cdot$ bahasa}
Kemajuan cepat pada pemrosesan bahasa alami. Transformer terbit 2017 -- bab 15 membongkarnya.
\end{icardaccent}
\end{minipage}%
\hfill
\begin{minipage}[t]{0.3253\textwidth}
\begin{icardaccent}
\cardh{2022 -- kini $\cdot$ generatif}
Gelombang aplikasi AI generatif yang mengubah cara orang bekerja sehari-hari.
\end{icardaccent}
\end{minipage}%
{\fontsize{9.4}{12.6}\selectfont\color{fg2}Urutan ini penting untuk perencanaan: kemampuan yang matang lebih dulu -- persepsi -- adalah yang \textbf{paling murah dan paling andal} dipakai hari ini. Yang paling baru justru yang paling mahal dan paling belum stabil.\par}\vskip 4pt
\begin{iband}[amber]
Di lapangan, banyak kasus pemakaian yang benar-benar mendesak sebetulnya jatuh di gelombang \textbf{pertama dan kedua} -- OCR dokumen, klasifikasi keluhan, deteksi anomali. \hilite{Tidak semuanya menuntut model generatif.}
\end{iband}
\end{frame}
\note{Slide ini sering mengubah arah diskusi anggaran. Tanyakan: kasus Anda sebetulnya butuh gelombang yang mana?}

\framekicker{Bagian 1.4 -- 1.5}
\begin{frame}[shrink=25]{Waspadai gembar-gembor jangka pendek}
\begin{minipage}[t]{0.485\textwidth}
{\fontsize{9.4}{12.6}\selectfont\color{fg2}\textbf{Ciri gembar-gembor sekarang}\par}\vskip 4pt
\begin{ibullets}
\item Ramalan pengangguran massal dan lonjakan produktivitas \textbf{10x-100x} -- yang belum terjadi.
\item Narasi AGI dan superintelijensi yang menyeret jadwal jadi tidak masuk akal.
\item Investasi AI -- terutama pusat data dan GPU -- melampaui \textbf{\$200 miliar per tahun}, sementara pendapatannya jauh tertinggal di kisaran \textbf{\$30 miliar}.
\end{ibullets}

\end{minipage}%
\hfill
\begin{minipage}[t]{0.485\textwidth}
{\fontsize{9.4}{12.6}\selectfont\color{fg2}\textbf{Dua musim dingin sebelumnya}\par}\vskip 4pt
\begin{ibullets}
\item \textbf{Pertama} -- Minsky, 1967: \textit{"Dalam satu generasi... persoalan menciptakan kecerdasan buatan pada dasarnya akan terpecahkan."} Pendanaan runtuh pada 1970-an.
\item \textbf{Kedua} -- \textit{expert system} pada 1980-an; sekitar 1985 perusahaan membelanjakan \textbf{lebih dari \$1 miliar setahun}. Pada awal 1990-an terbukti mahal dirawat, sulit diskalakan, dan sempit cakupannya.
\item \textbf{Sekarang} -- masih di fase optimisme. Menurut Chollet, kemunduran sebesar 1990-an \hilite{tidak mungkin terulang} -- AI sudah membuktikan nilainya. Kalau ada musim dingin, mestinya sangat ringan.
\end{ibullets}

\end{minipage}%
\begin{iquote}
\quotetext{Today's "artificial intelligence" is more accurately described as "cognitive automation." AI excels at solving problems with narrowly defined requirements.}\quotecite{Chollet \& Watson, bab 1}
\end{iquote}
\end{frame}
\note{Ini slide yang paling berguna saat peserta duduk di rapat anggaran. Pesannya bukan 'jangan investasi', melainkan 'taksir sesuai apa yang terbukti, bukan sesuai apa yang dijanjikan'.}

\framekicker{Bagian 1.10}
\begin{frame}[shrink=25]{Urutannya terbalik dari yang dikira orang}
\leadpar{Mudah mengira bahwa keberhasilan praktis AI generatif-lah yang melahirkan keyakinan akan AGI dalam waktu dekat. Menurut Chollet, \hilite{yang terjadi justru sebaliknya}}
\begin{isteps}
\item \textbf{2013} -- di kalangan elite teknologi sudah muncul kekhawatiran bahwa AGI akan tiba dalam beberapa tahun. Saat itu yang dianggap di jalur menuju ke sana adalah \textbf{DeepMind}, startup riset AI London yang kemudian diakuisisi Google.
\item \textbf{2015} -- keyakinan itulah yang mendorong berdirinya \textbf{OpenAI}, yang semula bermaksud menjadi penyeimbang sumber terbuka bagi DeepMind.
\item \textbf{2016} -- ajakan rekrutmen OpenAI menjanjikan \textbf{AGI tercapai pada 2020}. Adil untuk dicatat: hanya sebagian kecil orang di industri yang percaya jadwal seoptimistis itu waktu itu.
\item \textbf{Awal 2023} -- sebagian besar insinyur di San Francisco Bay Area tampak yakin AGI akan tiba dalam satu-dua tahun berikutnya.
\end{isteps}
\begin{iband}[amber]
OpenAI berperan penting menyalakan AI generatif. Jadi dalam pelintiran yang ganjil, \textbf{keyakinan akan AGI-lah yang menaikkan AI generatif}, bukan sebaliknya.
\end{iband}
\end{frame}
\note{Ini konteks yang jarang diketahui peserta dan sangat menolong saat membaca berita AI: banyak klaim yang beredar adalah keturunan dari keyakinan 2013, bukan simpulan dari bukti 2025.}

\framekicker{Bagian 1.5}
\begin{frame}[shrink=25]{Otomasi kognitif bukan kecerdasan}
\begin{iquote}
\quotetext{Intelligence is the ability to face the unknown, adapt to it, and learn from it. Automation, even at its best, can only handle situations it's been trained on.}\quotecite{Chollet \& Watson, bab 1}
\end{iquote}
\begin{itable}{>{\raggedright\arraybackslash}p{0.1692\textwidth}>{\raggedright\arraybackslash}p{0.3854\textwidth}>{\raggedright\arraybackslash}p{0.3854\textwidth}}
\thead{} & \thead{Otomasi kognitif (yang kita punya)} & \thead{Kecerdasan (yang belum)} \\ \midrule
Cakupan & Persoalan dengan syarat yang sempit dan jelas & Persoalan yang belum pernah dirumuskan siapa pun \\
Sumber kemampuan & Contoh dalam data pelatihan & Penyesuaian saat berhadapan dengan yang tak dikenal \\
Saat menemui hal baru & Menurun diam-diam, sering tanpa memberi tanda & Belajar dari situasi itu lalu memperbaiki diri \\
\end{itable}
\begin{minipage}[t]{0.485\textwidth}
{\fontsize{9.4}{12.6}\selectfont\color{fg2}Metafora Chollet: \textbf{AI itu seperti tokoh kartun, kecerdasan seperti makhluk hidup.} Kartun, sebagus apa pun gambarnya, hanya bisa memainkan adegan yang untuknya ia digambar. Makhluk hidup bisa menyesuaikan diri dengan yang tak terduga.\par}\vskip 4pt

\end{minipage}%
\hfill
\begin{minipage}[t]{0.485\textwidth}
{\fontsize{9.4}{12.6}\selectfont\color{fg2}\textit{"Kalau kartunnya digambar cukup realistis dan mencakup cukup banyak adegan, apa bedanya?"} Bedanya adalah \textbf{keluwesan menyesuaikan diri} -- dan itulah sebabnya membangun otomasi yang kukuh begitu sulit: ia menuntut Anda memperhitungkan \hilite{setiap skenario yang mungkin}\par}\vskip 4pt

\end{minipage}%
\begin{iband}[rose]
Pembedaan ini punya akibat langsung di produksi: kalau sistem hanya sanggup menangani apa yang pernah dilatihkan, maka \hilite{pemantauan pergeseran data bukan fitur tambahan}, melainkan syarat agar sistem tetap layak dipakai. Bab 18 kembali ke sini.
\end{iband}
{\fontsize{9.4}{12.6}\selectfont\color{fg2}Karena itu Chollet menutup bagian ini dengan tenang: jangan cemas AI tiba-tiba sadar diri lalu mengambil alih kemanusiaan. Teknologi hari ini tidak menuju ke sana. \textit{"Seperti mengharapkan jam yang lebih baik akan melahirkan perjalanan waktu -- keduanya hal yang sama sekali berbeda."}\par}\vskip 4pt
\end{frame}
\note{Ini jembatan ke seluruh sisa kursus. Kalimat kuncinya: model yang bagus di laboratorium bisa jadi berbahaya di produksi, dan bedanya ada pada apa yang tidak pernah ia lihat.}

\framekicker{Bagian 1.6}
\begin{frame}[shrink=25]{Janji jangka panjang}
{\fontsize{9.4}{12.6}\selectfont\color{fg2}Chollet menutup bab ini dengan ramalannya sendiri dari 2017, yang sebagian besar sudah terpenuhi pada 2025 -- dipakai bukan untuk menyombong, melainkan untuk menunjukkan bentuk taksiran yang benar: \textbf{arahnya tepat, jadwalnya yang meleset}.\par}\vskip 4pt
\begin{itable}{>{\raggedright\arraybackslash}p{0.376\textwidth}>{\raggedright\arraybackslash}p{0.564\textwidth}}
\thead{Ramalan 2017} & \thead{Keadaannya pada 2025} \\ \midrule
AI jadi asisten, bahkan teman & \textbf{Puluhan juta orang} memakai ChatGPT, Gemini, Claude sebagai asisten harian. Ratusan ribu berinteraksi dengan "teman" AI di aplikasi seperti Character.ai. \\
Menjawab pertanyaan, membantu mendidik anak & Ternyata \textbf{penjawaban pertanyaan dan bantuan pekerjaan rumah} justru menjadi pemakaian nomor satu chatbot ini. \\
Mengantar Anda dari A ke B & Kemudi otonom penuh sudah terpasang dalam skala nyata di Phoenix, San Francisco, Los Angeles, dan Austin. \\
Membantu ilmuwan menemukan terobosan & \textbf{AlphaFold} membantu biolog meramal struktur protein. Matematikawan \textbf{Terence Tao} memperkirakan AI bisa menjadi ko-penulis yang andal dalam riset matematika sekitar 2026. \\
\end{itable}
\begin{iband}[signal]
Kesimpulan bab: \hilite{gembar-gembor jangka pendek akan kempis, perubahan jangka panjang tetap datang}. Dua-duanya benar sekaligus, dan itulah yang membuat perencanaan jadi sulit.
\end{iband}
\end{frame}
\note{Tutup dengan mengingatkan bahwa sisa buku ini praktis semua. Bab 1 adalah satu-satunya bab yang boleh berdebat; setelah ini semuanya kode.}

\framekicker{Ringkasan}
\begin{frame}[shrink=25]{Yang wajib terbawa dari bab 1}
\begin{isteps}
\item \textbf{AI $\supset$ machine learning $\supset$ deep learning.} Symbolic AI ada di dalam AI, di luar ML.
\item \textbf{Machine learning membalik arah pemrograman}: data + jawaban masuk, aturan keluar.
\item \textbf{Tiga bahan wajib}: masukan, contoh keluaran, dan ukuran mutu. Yang paling sering hilang adalah yang ketiga.
\item \textbf{'Deep' = banyak lapis representasi}, bukan pemahaman yang lebih dalam.
\item \textbf{Bobot, fungsi rugi, backpropagation} -- lingkar pelatihan yang akan dibongkar bab 2.
\item \textbf{Capaiannya nyata, taksirannya sering meleset.} Ini otomasi kognitif, belum kecerdasan.
\end{isteps}
\vskip 4pt
\reslink{BAB BERIKUT}{Bab 2 --- Blok bangunan matematis}{../ch02/index.html}
\reslink{TEKS PENUH}{deeplearningwithpython.io}{https://deeplearningwithpython.io/chapters/chapter01\_what-is-deep-learning}
\reslink{KODE BUKU}{fchollet/deep-learning-with-python-notebooks}{https://github.com/fchollet/deep-learning-with-python-notebooks}
\vskip 2pt
\end{frame}
\note{Kalau waktu habis, enam butir ini yang harus sempat dibacakan.}

\end{document}
